% ==============================================================================
% 00_FRONTMATTER.TEX - PHẦN MỞ ĐẦU & QUY CHUẨN
% ==============================================================================

% Đánh số trang kiểu La Mã thường cho phần mở đầu
\pagenumbering{roman}
\setcounter{page}{1}

% ------------------------------------------------------------------------------
% 1. TRANG BÌA CHÍNH (BÌA NGOÀI)
% ------------------------------------------------------------------------------
\begin{titlepage}
\begin{center}
    % Khung viền đôi trang bìa chuẩn đại học Việt Nam
    \vspace*{-10mm}
    {\fontsize{13pt}{15pt}\selectfont \textbf{BỘ GIAO THÔNG VẬN TẢI}}\\[3mm]
    {\fontsize{14pt}{16pt}\selectfont \textbf{TRƯỜNG ĐẠI HỌC GIAO THÔNG VẬN TẢI}}\\[2mm]
    {\fontsize{13pt}{15pt}\selectfont \textbf{KHOA CÔNG NGHỆ THÔNG TIN}}\\[1mm]
    {\fontsize{12pt}{14pt}\selectfont \textbf{BỘ MÔN CÔNG NGHỆ PHẦN MỀM}}\\[6mm]

    \includegraphics[width=35mm, height=35mm, keepaspectratio]{images/logo_utc.png}\\[8mm]

    {\fontsize{18pt}{22pt}\selectfont \textbf{BÁO CÁO BÀI TẬP LỚN}}\\[4mm]
    {\fontsize{14pt}{17pt}\selectfont \textbf{MÔN HỌC: PHÁT TRIỂN PHẦN MỀM MÃ NGUỒN MỞ}}\\[8mm]

    \noindent\rule{\textwidth}{1.5pt}\\[4mm]
    {\fontsize{16pt}{20pt}\selectfont \textbf{ĐỀ TÀI:}}\\[3mm]
    {\fontsize{18pt}{22pt}\selectfont \textbf{XÂY DỰNG WEBSITE POS QUẢN LÝ BÁN HÀNG TRÊN NỀN TẢNG SPRING BOOT VÀ REACTJS}}\\[3mm]
    \noindent\rule{\textwidth}{1.5pt}\\[15mm]

    \begin{minipage}{0.9\textwidth}
        \fontsize{13pt}{16pt}\selectfont
        \begin{tabular}{p{5.5cm}l}
            \textbf{Giảng viên hướng dẫn:} & \textbf{ThS. Đào Thị Lệ Thủy} \\[3mm]
            \textbf{Sinh viên thực hiện 1:} & \dots\dots\dots\dots\dots\dots\dots\dots\dots\dots\dots\dots\dots\dots \\[2mm]
            \textbf{Mã số sinh viên 1:}     & \dots\dots\dots\dots\dots\dots\dots\dots\dots\dots\dots\dots\dots\dots \\[2mm]
            \textbf{Sinh viên thực hiện 2:} & \dots\dots\dots\dots\dots\dots\dots\dots\dots\dots\dots\dots\dots\dots \\[2mm]
            \textbf{Mã số sinh viên 2:}     & \dots\dots\dots\dots\dots\dots\dots\dots\dots\dots\dots\dots\dots\dots \\[2mm]
            \textbf{Lớp:}                  & \dots\dots\dots\dots\dots\dots\dots\dots\dots\dots\dots\dots\dots\dots \\
        \end{tabular}
    \end{minipage}

    \vfill
    {\fontsize{13pt}{15pt}\selectfont \textbf{HÀ NỘI, NĂM HỌC 2026 - 2027}}
\end{center}
\end{titlepage}

% ------------------------------------------------------------------------------
% 2. TRANG BÌA PHỤ (BÌA TRONG)
% ------------------------------------------------------------------------------
\cleardoublepage
\begin{titlepage}
\begin{center}
    \vspace*{-10mm}
    {\fontsize{13pt}{15pt}\selectfont \textbf{BỘ GIAO THÔNG VẬN TẢI}}\\[3mm]
    {\fontsize{14pt}{16pt}\selectfont \textbf{TRƯỜNG ĐẠI HỌC GIAO THÔNG VẬN TẢI}}\\[2mm]
    {\fontsize{13pt}{15pt}\selectfont \textbf{KHOA CÔNG NGHỆ THÔNG TIN}}\\[15mm]

    {\fontsize{18pt}{22pt}\selectfont \textbf{BÁO CÁO BÀI TẬP LỚN}}\\[4mm]
    {\fontsize{14pt}{17pt}\selectfont \textbf{MÔN HỌC: PHÁT TRIỂN PHẦN MỀM MÃ NGUỒN MỞ}}\\[10mm]

    \noindent\rule{\textwidth}{1pt}\\[4mm]
    {\fontsize{18pt}{22pt}\selectfont \textbf{XÂY DỰNG WEBSITE POS QUẢN LÝ BÁN HÀNG TRÊN NỀN TẢNG SPRING BOOT VÀ REACTJS}}\\[4mm]
    \noindent\rule{\textwidth}{1pt}\\[25mm]

    \begin{minipage}{0.9\textwidth}
        \fontsize{13pt}{16pt}\selectfont
        \begin{tabular}{p{5.5cm}l}
            \textbf{Giảng viên hướng dẫn:} & \textbf{ThS. Đào Thị Lệ Thủy} \\[4mm]
            \textbf{Sinh viên thực hiện:}   & \\[2mm]
            & 1. \dots\dots\dots\dots\dots\dots\dots\dots\dots\dots - MSSV: \dots\dots\dots\dots \\[2mm]
            & 2. \dots\dots\dots\dots\dots\dots\dots\dots\dots\dots - MSSV: \dots\dots\dots\dots \\[3mm]
            \textbf{Khóa:}                  & K64 \\[2mm]
            \textbf{Hệ đào tạo:}           & Đại học chính quy \\
        \end{tabular}
    \end{minipage}

    \vfill
    {\fontsize{13pt}{15pt}\selectfont \textbf{HÀ NỘI, THÁNG 09/2026}}
\end{center}
\end{titlepage}

% ------------------------------------------------------------------------------
% 3. LỜI CAM ĐOAN
% ------------------------------------------------------------------------------
\cleardoublepage
\chapter*{LỜI CAM ĐOAN}
\addcontentsline{toc}{chapter}{LỜI CAM ĐOAN}

Chúng tôi xin cam đoan rằng:
\begin{enumerate}[label=\arabic*.]
    \item Báo cáo bài tập lớn đề tài \textit{''Xây dựng website POS Quản lý bán hàng trên nền tảng Spring Boot và ReactJS''} là công trình nghiên cứu và thực hiện độc lập của nhóm chúng tôi dưới sự hướng dẫn chuyên môn của cô \textbf{ThS. Đào Thị Lệ Thủy}.
    \item Toàn bộ số liệu khảo sát, sơ đồ kiến trúc, từ điển dữ liệu và kết quả thực nghiệm được trình bày trong báo cáo là trung thực, phản ánh chính xác mã nguồn chương trình thực tế của dự án.
    \item Các đoạn mã nguồn tham khảo từ các dự án mã nguồn mở, tài liệu kỹ thuật, thư viện bên thứ ba (Spring Framework, ReactJS, Vite, Tailwind CSS, PayOS, AWS SDK) đều đã được trích dẫn nguồn gốc rõ ràng tại danh mục Tài liệu tham khảo.
    \item Chúng tôi hoàn toàn chịu trách nhiệm trước Nhà trường và Khoa Công nghệ Thông tin về tính xác thực của các nội dung được cam đoan trên.
\end{enumerate}

\vspace{15mm}
\begin{flushright}
    \textit{Hà Nội, ngày \dots\dots tháng 09 năm 2026}\\[3mm]
    \textbf{Đại diện nhóm sinh viên thực hiện}\\[15mm]
    \textit{(Ký và ghi rõ họ tên)}
\end{flushright}

% ------------------------------------------------------------------------------
% 4. LỜI CẢM ƠN
% ------------------------------------------------------------------------------
\cleardoublepage
\chapter*{LỜI CẢM ƠN}
\addcontentsline{toc}{chapter}{LỜI CẢM ƠN}

Lời đầu tiên, nhóm chúng em xin gửi lời cảm ơn chân thành và sâu sắc nhất tới cô \textbf{ThS. Đào Thị Lệ Thủy} – giảng viên phụ trách môn học \textit{Phát triển phần mềm mã nguồn mở}. Trong suốt quá trình học tập và hoàn thiện đề tài bài tập lớn, cô đã luôn nhiệt tình định hướng chuyên môn, gợi mở các giải pháp kỹ thuật hiện đại và đưa ra những lời nhận xét, góp ý quý báu giúp hệ thống của chúng em được hoàn thiện toàn diện cả về mặt kiến trúc lẫn trải nghiệm người dùng thực tế.

Chúng em cũng xin trân trọng cảm ơn các Thầy, Cô giáo trong \textbf{Khoa Công nghệ Thông tin – Trường Đại học Giao thông Vận tải} đã tận tâm truyền đạt nền tảng kiến thức vững chắc về công nghệ phần mềm, lập trình hướng đối tượng, phân tích thiết kế hệ thống và cơ sở dữ liệu trong suốt các năm học vừa qua.

Mặc dù nhóm đã hết sức nỗ lực nghiên cứu mã nguồn, phân tích nghiệp vụ và xây dựng ứng dụng với tinh thần trách nhiệm cao nhất, song do hạn chế về thời gian cũng như kinh nghiệm thực tiễn, báo cáo và hệ thống chắc chắn khó tránh khỏi những thiếu sót nhất định. Nhóm chúng em rất mong nhận được những ý kiến đóng góp, chỉ bảo quý báu từ quý Thầy Cô để sản phẩm tiếp tục được nâng cấp và hoàn thiện hơn trong tương lai.

\vspace{5mm}
\noindent Chúng em xin kính chúc cô ThS. Đào Thị Lệ Thủy cùng toàn thể các Thầy Cô giáo trong Khoa luôn dồi dào sức khỏe, hạnh phúc và gặt hái được nhiều thành công hơn nữa trong sự nghiệp trồng người!

\vspace{15mm}
\begin{flushright}
    \textbf{Tập thể nhóm sinh viên thực hiện}
\end{flushright}

% ------------------------------------------------------------------------------
% 5. BẢNG PHÂN CÔNG NHIỆM VỤ THÀNH VIÊN
% ------------------------------------------------------------------------------
\cleardoublepage
\chapter*{BẢNG PHÂN CÔNG NHIỆM VỤ THÀNH VIÊN}
\addcontentsline{toc}{chapter}{BẢNG PHÂN CÔNG NHIỆM VỤ THÀNH VIÊN}

\begin{table}[htbp]
\centering
\caption{Bảng phân công nhiệm vụ và đánh giá mức độ đóng góp}
\label{tab:phan_cong_nhiem_vu}
\vspace{2mm}
\begin{tabularx}{\textwidth}{|c|p{3.5cm}|c|X|c|}
\hline
\textbf{STT} & \textbf{Họ và tên} & \textbf{MSSV} & \textbf{Nhiệm vụ đảm nhiệm chính} & \textbf{Tỷ lệ (\%)} \\ \hline
1 & \dots\dots\dots\dots\dots\dots & \dots\dots\dots\dots & 
\begin{itemize}[leftmargin=*, nosep]
    \item Thiết kế kiến trúc tổng thể Backend (Java 25, Spring Boot 4.1.0).
    \item Thiết kế và chuẩn hóa CSDL MySQL, triển khai Flyway migration.
    \item Cài đặt cơ chế bảo mật JWT + Refresh Token Rotation, kiểm toán Activity Logs.
    \item Tích hợp cổng thanh toán VietQR PayOS và lưu trữ file AWS S3.
    \item Soạn thảo Chương 1 và Chương 2 của báo cáo.
\end{itemize} & 100\% \\ \hline
2 & \dots\dots\dots\dots\dots\dots & \dots\dots\dots\dots & 
\begin{itemize}[leftmargin=*, nosep]
    \item Thiết kế và xây dựng giao diện người dùng Frontend (React 19, Vite, Tailwind CSS v4).
    \item Triển khai quản lý Server State qua TanStack Query v5, polling giao dịch QR.
    \item Tối ưu hóa giao diện bán hàng nhanh POS, luồng in hóa đơn nhiệt 80mm.
    \item Thiết kế kịch bản kiểm thử (Test Cases Matrix) và thực nghiệm hệ thống.
    \item Soạn thảo Chương 3, Kết luận và tổng hợp báo cáo.
\end{itemize} & 100\% \\ \hline
\end{tabularx}
\end{table}

\begin{flushright}
    \textit{Xác nhận của nhóm trưởng: Toàn bộ thành viên đều tham gia đầy đủ và đạt kết quả tốt.}
\end{flushright}

% ------------------------------------------------------------------------------
% 6. DANH MỤC THUẬT NGỮ VÀ TỪ VIẾT TẮT
% ------------------------------------------------------------------------------
\cleardoublepage
\chapter*{DANH MỤC TỪ VIẾT TẮT}
\addcontentsline{toc}{chapter}{DANH MỤC TỪ VIẾT TẮT}

\begin{longtable}{|c|p{3.5cm}|p{9cm}|}
\caption{Danh mục các từ viết tắt sử dụng trong báo cáo} \label{tab:viet_tat} \\
\hline
\textbf{Từ viết tắt} & \textbf{Tiếng Anh} & \textbf{Giải thích nghĩa tiếng Việt} \\ \hline
\endfirsthead
\hline
\textbf{Từ viết tắt} & \textbf{Tiếng Anh} & \textbf{Giải thích nghĩa tiếng Việt} \\ \hline
\endhead
\hline
\endfoot
\hline
\endlastfoot

ACID & Atomicity, Consistency, Isolation, Durability & Bốn thuộc tính đảm bảo tính tin cậy của giao dịch cơ sở dữ liệu. \\ \hline
ADR & Architecture Decision Record & Bản ghi chép các quyết định kiến trúc phần mềm cốt lõi. \\ \hline
API & Application Programming Interface & Giao diện lập trình ứng dụng. \\ \hline
AWS & Amazon Web Services & Nền tảng dịch vụ điện toán đám mây của Amazon. \\ \hline
BFD & Business Function Diagram & Biểu đồ phân rã chức năng nghiệp vụ. \\ \hline
BOGO & Buy One, Get One & Chương trình khuyến mãi mua sản phẩm A tặng/giảm giá sản phẩm B. \\ \hline
CRM & Customer Relationship Management & Hệ thống quản lý quan hệ khách hàng và tích lũy điểm thưởng. \\ \hline
CSDL / DBMS & Database Management System & Hệ quản trị cơ sở dữ liệu. \\ \hline
DTO & Data Transfer Object & Đối tượng đóng gói truyền tải dữ liệu giữa Client và Server. \\ \hline
ERD & Entity Relationship Diagram & Sơ đồ mô hình quan hệ thực thể. \\ \hline
HttpOnly & HTTP-Only Cookie & Thuộc tính bảo mật cookie ngăn chặn truy cập mã độc qua JavaScript/XSS. \\ \hline
JPA & Jakarta Persistence API & Chuẩn giao tiếp ánh xạ đối tượng vào bảng CSDL (ORM) trong Java. \\ \hline
JSON & JavaScript Object Notation & Định dạng trao đổi dữ liệu dạng văn bản nhẹ. \\ \hline
JWT & JSON Web Token & Chuẩn mã hóa chuỗi ký điện tử dùng để xác thực người dùng. \\ \hline
MVC & Model - View - Controller & Mô hình kiến trúc phân tách Dữ liệu - Giao diện - Điều khiển. \\ \hline
PayOS & PayOS Payment Gateway & Cổng thanh toán trực tuyến hỗ trợ tạo mã VietQR động theo chuẩn NAPAS. \\ \hline
POS & Point of Sale & Điểm bán lẻ, phần mềm quản lý thu ngân và xuất bán hàng hóa tại quầy. \\ \hline
RBAC & Role-Based Access Control & Cơ chế kiểm soát và phân quyền truy cập dựa trên vai trò người dùng. \\ \hline
REST & Representational State Transfer & Kiểu kiến trúc tiêu chuẩn cho các dịch vụ Web phân tán qua HTTP. \\ \hline
S3 & Simple Storage Service & Dịch vụ lưu trữ đối tượng đám mây có độ bền cao của Amazon. \\ \hline
SKU & Stock Keeping Unit & Đơn vị phân loại hàng hóa tồn kho dựa trên mã biến thể duy nhất. \\ \hline
UML & Unified Modeling Language & Ngôn ngữ mô hình hóa thống nhất phục vụ phân tích thiết kế phần mềm. \\ \hline
URI / URL & Uniform Resource Identifier / Locator & Định danh và địa chỉ định vị tài nguyên mạng. \\ \hline
VAT & Value Added Tax & Thuế giá trị gia tăng được áp dụng trên hóa đơn (10\%). \\ \hline
VietQR & Vietnam Quick Response code & Tiêu chuẩn nhận diện mã thanh toán QR quốc gia do NAPAS ban hành. \\
\end{longtable}

% ------------------------------------------------------------------------------
% 7. MỤC LỤC, DANH MỤC HÌNH ẢNH & BẢNG BIỂU TỰ ĐỘNG
% ------------------------------------------------------------------------------
\cleardoublepage
\tableofcontents

\cleardoublepage
\listoffigures
\addcontentsline{toc}{chapter}{DANH MỤC HÌNH ẢNH}

\cleardoublepage
\listoftables
\addcontentsline{toc}{chapter}{DANH MỤC BẢNG BIỂU}
